\section{Integral coefficient lattices and the ramified conductor}
\label{sec:integral-coefficient-conductor}

The following maps retain the arithmetic of the source-displayed coefficient
\eqref{eq:gct-nrh-coefficient-original}. Put
\[
 O=\mathbb Z[1/2],\quad R=O[q,q^{-1}],\quad t=q-1,\quad
 r=q-q^{-1},\quad h=[7]_q[3]_q,\quad C=r^2h.
\]
The letters $r$ and $a$ in this section use the original $b=c=0$
root coordinate and base map $a=-r^2$, respectively. The trace coordinate
is still $p=q+q^{-1}$. None of $3,7,h$ is inverted in $R$.

\subsection{The exact lattice map over the integral base}

The two-sheet calculation of Section~\ref{sec:specialization} holds over
$O[a]$: the displayed monic bases, the component restriction maps, and
their inverses only use division by $2$. Base change by
$O[a]\to R$, $a\mapsto-r^2$, gives free modules $B_R,N_R$ with their
original ordered bases and map
\[
 V:B_R\longrightarrow N_R,\qquad [V]=\operatorname{diag}(1,1,1,r^2).
\]
This assertion is also checked directly: the last component vector
$r_1r_2$ maps to $-a(1,-1)=r^2(1,-1)$. Thus the original sign is retained.
Since $R$ is a domain and $r\ne0$, this map is injective, and projecting
onto the last coordinate modulo $r^2$ proves
\[
 0\longrightarrow B_R\xrightarrow{V}N_R
 \longrightarrow R/(r^2)\longrightarrow0.
\]

Let $E_C=R^4$ be a free module with ordered basis $e_1,e_2,e_3,e_4$.
Define all entries of the following maps:
\[
 H:E_C\longrightarrow B_R,\quad[H]=\operatorname{diag}(1,1,1,h),
 \qquad D_C:E_C\longrightarrow N_R,\quad[D_C]=\operatorname{diag}(1,1,1,C).
\]
Matrix multiplication proves $D_C=VH$; both maps are injective because
$h,C$ are nonzero. Define the two actual image lattices
$L_C=D_C(E_C)$ and $L_r=V(B_R)$ inside $N_R$.
Their first three coordinates equal $R$ and their last coordinates are
$CR$ and $r^2R$, respectively. Consequently there is an exact sequence
\begin{equation}\label{eq:integral-lattice-sequence}
 0\longrightarrow R/(h)\xrightarrow{\iota}R/(C)
 \xrightarrow{\pi}R/(r^2)\longrightarrow0,
 \quad \iota(\bar f)=\overline{r^2f},\quad\pi(\bar f)=\bar f.
\end{equation}
Here $R/(h)=L_r/L_C$, $R/(C)=N_R/L_C$ and $R/(r^2)=N_R/L_r$ under
the indicated last-coordinate maps. For well-definedness of $\iota$,
replacing $f$ by $f+hg$ changes its image by $Cg$. If $r^2f=Cg$,
cancellation of the nonzero $r^2$ gives $f=hg$, proving injectivity.
The kernel of $\pi$ consists exactly of classes $r^2f$, and $\pi$ is
surjective by lifting a representative. This proves every assertion.

The sign action $\Gamma=\mu_2\times\mu_2$ on the four original basis
vectors has characters $1,\chi_1,\chi_2,\chi_1\chi_2$.
Let it act trivially on $R$ and by these same characters on $E_C$.
All diagonal maps above commute with it, so every last-coordinate
quotient in \eqref{eq:integral-lattice-sequence} has character
$\chi_1\chi_2$. The semilinear involution $q\mapsto q^{-1}$, fixing
the four chosen basis vectors, also commutes with the maps: it sends
$r$ to $-r$ and fixes $r^2,h,C$. These are explicit symmetries of this
lattice diagram. No unspecified Chevalley action on $E_C$ is inferred.

\subsection{The full specialization sequence and its factor 21}

Specialization means tensoring over $R$ with $O=R/(t)$.
Write $u=(q+1)/q$, so $r=tu$ and $u(1)=2$.
We have $h(1)=21$. Both $V$ and $D_C$ specialize to
$\operatorname{diag}(1,1,1,0)$. Their kernels and cokernels are therefore
free copies of $O$ with basis given by the fourth coordinate.
The induced kernel map from $H$ is multiplication by $21$; the induced
cokernel map from the identity on $N_R$ is the identity of $O$.
In particular this diagram preserves more than the ranks of the maps.

Here is the complete derived specialization of
\eqref{eq:integral-lattice-sequence}, including the connecting map.
For any $R$-module $M$, the two-term free resolution
$0\to R\xrightarrow{t}R\to O\to0$ proves
\[
 \operatorname{Tor}_1^R(M,O)=\{m\in M:tm=0\},\qquad
 \operatorname{Tor}_j^R(M,O)=0\quad(j\ge2).
\]
If $d=td_1\ne0$, cancellation in $R$ shows
\[
 \operatorname{Tor}_1^R(R/(d),O)=O\,\overline{d_1}.
\]
Indeed $tf=dg=td_1g$ implies $f=d_1g$, and $d_1g\in(d)$ holds
exactly when $g\in(t)$. The ideal $(t)$ is prime because $R/(t)=O$
is a domain. It does not contain $h$, whose value is $21\ne0$ in $O$.
Thus $tf=hg$ forces $g=tg_1$, then $f=hg_1$, and
$\operatorname{Tor}_1^R(R/(h),O)=0$. Applying these facts to $d=C,r^2$ proves the
exact specialized sequence
\begin{equation}\label{eq:integral-specialization-21}
 0\longrightarrow O\xrightarrow{\,21\,}O
 \xrightarrow{\mathrm{mod}\ 21}O/(21)
 \xrightarrow{\,0\,}O\xrightarrow{\,1\,}O\longrightarrow0.
\end{equation}
We verify the maps, so that no exactness assertion depends on a hidden
choice of generators. The two Tor generators are $C/t=tu^2h$ and
$r^2/t=tu^2$, respectively. Projection sends the first to $h$ times the
second, and hence gives $21$ after reduction by $t$.
To compute the connecting map, lift $r^2/t$ from $R/(r^2)$ to
$R/(C)$. Its product by $t$ is $r^2=\iota(1)$; the connecting class
is therefore $1$ in $R/(h,t)=O/(21)$, with positive sign. The next map
is induced by multiplication by $r^2$, which vanishes modulo $t$.
The final map is the identity on constant classes. Multiplication by
$21$ is injective on $O$, with cokernel exactly $O/(21)$, verifying
exactness of the displayed sequence directly as well.

Over $D=\mathbb Q[q,q^{-1}]_{(q-1)}$ the element $h$ is a unit: its value
$21$ is nonzero in the residue field $\mathbb Q$. Thus $H$ has the explicit
inverse $\operatorname{diag}(1,1,1,h^{-1})$ there, and it identifies
the two injections over $D$. Over $R$ the additional quotient $R/(h)$
and the residual $O/(21)$ remain. These are the additional comparison
modules killed by passage to $D$. This is not an exhaustive inventory
of localization effects: the same localization also removes the $q=-1$
branch of $r^2$.

For completeness the entire $t$-saturation can be computed while retaining
the integral constants. Work in
$R_+=O[q,q^{-1},(q+1)^{-1}]$, where $u$ is a unit and $(t)$ is still
prime. The ideal generated by $C$ is $(t^2h)$ and the ideal generated
by $r^2$ is $(t^2)$. Their $t$-saturations are $(h)$ and $R_+$.
For the first assertion, if $t^ef=t^2hg$, then for $e\ge2$ cancellation
gives $t^{e-2}f=hg$. Primality of $t$ and $t\nmid h$ imply successively
that $g$ is divisible by $t^{e-2}$, so $f\in(h)$. For $e<2$ the same
conclusion follows directly by cancellation. Conversely $t^2hR_+$ lies
in $(C)$, proving the asserted saturation; the second assertion follows
from $t^2R_+=(r^2)$. Thus the saturated image lattices are
$R_+^3\oplus hR_+$ and $R_+^4$, respectively.

\subsection{The retained integral jets and the two exceptional primes}

The exact equality $C=t^2u^2h$ gives
\[
 C\equiv84t^2\pmod{t^3R_+}.
\]
This follows by evaluating $u^2h$ at $q=1$: $2^2\cdot7\cdot3=84$.
Over $\mathbb Q$ the coefficient has exact order two, as previously proved.
The integral residue $84$ also determines why that order changes in
characteristics $3$ and $7$; these characteristics are allowed over $O$.

For an odd prime $\ell$, in $\mathbb F_\ell[q,q^{-1}]$ one has the polynomial
identity
\begin{equation}\label{eq:integral-quantum-prime}
 [\ell]_q=(q-q^{-1})^{\ell-1}.
\end{equation}
To prove it, multiply the finite sum defining $[\ell]_q$ by $q^{\ell-1}$.
The result is $1+q^2+\cdots+q^{2(\ell-1)}$.
Its product with $q^2-1$ is $q^{2\ell}-1=(q^2-1)^\ell$ in
characteristic $\ell$. Cancellation in the polynomial domain gives
$(q^2-1)^{\ell-1}$, and multiplication by $q^{-(\ell-1)}$ proves
\eqref{eq:integral-quantum-prime}. In particular the original coefficient
reduces exactly to
\[
 C=r^4[7]_q\quad\text{over }\mathbb F_3,\qquad
 C=r^8[3]_q\quad\text{over }\mathbb F_7.
\]
Because $r=tu$, $u(1)=2$, $[7]_1=1$ in $\mathbb F_3$ and $[3]_1=3$ in
$\mathbb F_7$, the exact orders of $C$ at $q=1$ are four and eight,
respectively. Their first nonzero coefficients are
$2^4\cdot1=1$ in $\mathbb F_3$ and $2^8\cdot3=5$ in $\mathbb F_7$.
For every other odd characteristic $\ell$, both $21$ and $4$ are
nonzero, so the exact order is two and its first coefficient is
$84\bmod\ell$. The node map still has last entry $r^2$, of exact
order two in every odd characteristic. The larger orders therefore
belong to the retained residual factor of the GCT source-displayed polynomial,
and the injection $H$ records their relation to the node map.

These results construct an integral coefficient-lattice correspondence.
They specify every map and every residual module in that correspondence.
Section~\ref{sec:gct-generator-interval} supplies the actual Cartan-twisted
divided-square endpoint from the original GCT-IV basis subquotient.
Adding its three specified identity summands gives exactly $D_C$ here,
so $D_C=VH$ is also an explicit comparison of that generator endpoint
with the original node. The trace obstruction proved there prevents
extending the four-dimensional interval to the full raising action.
Section~\ref{sec:full-source-generators} now constructs that full original
source module, all its raising and lowering operators, and the exact
integral quotient map to the retained interval. Its raw-word projection
has a specified section and free kernel; the square comparing its full
domain with the earlier restricted four-word lattice retains the latter's
nonzero cokernel. Thus the larger source action and the conductor
comparison here are related by proved integral maps.
